\section{第一阶段源码走读：x86-64 \#PF 到 VMA、PTE 与 COW}

前面已经讲过虚拟地址、VMA、PTE、TLB 和 COW，但只讲概念还不够。真正读懂操作系统，必须能把一次用户态 load/store 失败拆成硬件异常、入口现场、VMA 判定、PTE 状态、缺页修复、TLB 失效和返回重试。x86-64 上这条路径的核心异常是 \#PF，Linux 里它从 \filepath{arch/x86/mm/fault.c} 进入，再落到 \filepath{mm/memory.c} 的通用内存管理逻辑。

\begin{keyidea}
\#PF 不是“程序访问错了内存”的同义词。它只是 CPU 告诉内核：“这次地址翻译或权限检查没有按当前页表完成。”内核必须结合 CR2、error code、当前上下文、VMA 和 PTE，才能决定是分配零页、读文件页、执行 COW、返回 \code{-EFAULT}，还是给进程发送信号。
\end{keyidea}

\subsection{硬件交给内核的最小证据}

x86-64 上发生 page fault 时，CPU 至少提供三类证据：faulting virtual address、错误码和被打断现场。faulting address 由 CR2 保存；错误码说明这是不是 present 页上的保护错误、是不是写访问、是不是用户态访问、是不是取指访问等；入口保存的寄存器现场说明 fault 发生在什么 RIP、什么 RSP、什么 CPL。

\begin{longtable}{p{0.22\linewidth}p{0.34\linewidth}p{0.32\linewidth}}
\toprule
证据 & 典型来源 & OS 需要回答的问题 \\
\midrule
fault address & CR2 & 这个地址属于哪个 VMA，是否 canonical \\
error code.P & page fault error code & 是 PTE 不存在，还是 present 页权限拒绝 \\
error code.W/R & page fault error code & 是读、写还是取指触发 \\
error code.U/S & page fault error code & 是 Ring 3 访问，还是内核路径访问 \\
trap frame & 异常入口保存 & 返回后重试哪条指令，还是转成信号/错误 \\
\bottomrule
\end{longtable}

这张表解释了为什么“缺页处理”不能只写成 \code{if page missing then allocate}。同样是 \#PF，缺少 PTE 可能是匿名懒分配，present 页写保护可能是 COW，用户态访问无 VMA 是 \code{SIGSEGV}，内核 \code{copy\_from\_user} 访问坏地址应该返回 \code{-EFAULT}，内核普通指针 fault 则可能是 oops。

\subsection{源码地图：从 x86-64 异常到通用 mm}

阅读 Linux 时，先不要在整个源码树里搜索“page fault”到处跳。应该按层次看：

\begin{longtable}{p{0.28\linewidth}p{0.30\linewidth}p{0.30\linewidth}}
\toprule
路径 & 角色 & 阅读重点 \\
\midrule
\filepath{arch/x86/mm/fault.c} & x86-64 page fault 入口 & 读取 CR2、解释 error code、区分用户/内核上下文 \\
\filepath{arch/x86/entry/entry\_64.S} & 异常入口汇编 & trap frame、寄存器保存、进入 C 函数前的状态 \\
\filepath{include/linux/mm\_types.h} & mm 核心结构 & \code{mm\_struct}、\code{vm\_area\_struct}、页表相关对象 \\
\filepath{mm/mmap.c} & VMA 创建与查找 & \code{mmap}、\code{mprotect}、VMA 合并拆分 \\
\filepath{mm/memory.c} & 通用 fault 主干 & \code{handle\_mm\_fault}、匿名 fault、文件 fault、COW \\
\filepath{mm/rmap.c} & 反向映射 & 页被哪些映射引用，回收和 COW 需要的关系 \\
\filepath{mm/filemap.c} & 文件页 fault & page cache、文件映射、读入文件页 \\
\filepath{arch/x86/mm/tlb.c} & TLB 失效 & 本地 flush、远端 shootdown、PCID 相关路径 \\
\bottomrule
\end{longtable}

这条地图的分层很重要：\filepath{arch/x86/mm/fault.c} 负责架构相关入口，\filepath{mm/memory.c} 负责大部分架构无关语义。不要把二者混成一个函数理解。x86-64 入口知道 CR2 和 error code；通用 mm 知道 VMA、folio、匿名页、文件页、COW、回收和锁。

\subsection{第一步：地址形式和 VMA 查找不是同一件事}

用户态传来的值先要满足 x86-64 canonical address 规则。非 canonical 地址不是“找不到 VMA”那么简单，它在页表遍历前就可以被硬件拒绝。通过这一关后，内核才用 fault address 查找当前进程的 VMA。

\begin{lstlisting}
page fault classification:
  address is not canonical:
    reject before treating it as a normal VMA miss
  no VMA covers address:
    user instruction -> SIGSEGV
    kernel usercopy -> -EFAULT
  VMA covers address but permissions reject access:
    protection failure, not demand allocation
  VMA covers address and permissions allow access:
    inspect PTE and repair if possible
\end{lstlisting}

VMA 是虚拟地址区间的语义说明。它回答“这段地址是否应该存在，允许读写执行吗，来自匿名内存还是文件，私有还是共享”。PTE 是当前硬件翻译状态。一个合法 VMA 允许暂时没有 PTE；一个没有 VMA 的地址不能靠分配页修复。

\subsection{第二步：error code 决定走缺失路径还是保护路径}

page fault error code 的 present/protection 位把路径分成两类。若不是 present 页，常见情况是懒分配、文件页未读入或 swap 页未换入。若是 present 页上的保护 fault，常见情况是 COW、写只读映射、用户态访问 supervisor 页或取指违反 NX。

\begin{longtable}{p{0.20\linewidth}p{0.35\linewidth}p{0.33\linewidth}}
\toprule
硬件现象 & 内核可能动作 & 不能误判成什么 \\
\midrule
not-present 读匿名页 & 分配零页，安装 PTE，重试指令 & 程序错误 \\
not-present 读文件页 & 查 page cache，必要时发起 I/O & 简单 malloc \\
present 写只读 PTE & 若 VMA 允许写且 PTE 是 COW，则复制或提升 & 普通权限失败 \\
present 写只读 VMA & 发送信号或返回错误 & COW \\
present 取指 NX 页 & 发送信号或安全违规处理 & 数据页缺失 \\
\bottomrule
\end{longtable}

这里的关键是同时看 VMA 和 PTE。VMA 允许写、PTE 不可写，可能是 COW；VMA 本身不允许写，PTE 不可写就是权限拒绝。只看 PTE 会把只读文件映射误当成 COW；只看 VMA 会看不见硬件当前状态。

\subsection{匿名懒分配：合法地址为什么第一次访问才有物理页}

\code{mmap} 或堆增长可以先建立 VMA，不立刻分配物理页。第一次读写触发 not-present \#PF，内核分配一个清零页，安装 PTE，返回用户态重试 faulting instruction。用户程序感觉这条 load/store 成功了，但中间发生过一次异常和内核修复。

\begin{lstlisting}
anonymous demand fault:
  find VMA covering fault address
  verify VMA permits the access
  allocate zero-filled physical page
  install user PTE with VMA-derived permissions
  return to the same user RIP
  CPU retries the faulting load/store
\end{lstlisting}

这就是为什么 VIRT 可以很大而 RSS 很小。VMA 代表承诺，RSS 代表当前真的驻留的物理页。OS 用这个策略避免提前为大量可能永远不会访问的地址分配内存。

\subsection{文件映射缺页：read 和 mmap 在 page cache 汇合}

文件 \code{mmap} 后，用户第一次读某个页面也会触发 not-present \#PF。内核根据 VMA 找到文件和偏移，再到 page cache 查对应页。若 page cache 命中，这是 minor fault；若需要从存储设备读入，这是 major fault。读入后安装 PTE，用户指令重试。

\begin{lstlisting}
file-backed fault:
  fault address -> VMA
  VMA file + page offset -> page-cache index
  if page cache hit:
    map cached folio and return
  else:
    submit storage read, wait for completion
    mark folio uptodate
    map folio and return
\end{lstlisting}

因此 \code{read} 和 \code{mmap} 共享的不只是“文件内容”，而是内核里的 page cache 对象。前者通常把 page cache 内容复制到用户 buffer；后者把 page cache 页映射进用户地址空间。理解这一点，后面讲 page cache 回收、dirty、writeback 和文件系统日志才不会断开。

\subsection{COW 的完整条件：VMA 允许写，PTE 暂时禁止写}

fork 后父子共享匿名物理页。为了发现第一次写，父子 PTE 都要清掉 writable，并标记 COW。用户写入时，CPU 看到 present PTE 但不可写，于是产生保护型 \#PF。内核确认 VMA 允许写、PTE 是 COW，再决定复制还是提升。

\begin{lstlisting}
COW write fault:
  write to present but read-only PTE
  VMA says private mapping is writable
  PTE carries COW state
  if physical page refcount > 1:
    allocate new page and copy old content
    decrement old page refcount
    install writable PTE to new page
  else:
    only restore writable bit
  invalidate stale TLB translation
  return and retry the store
\end{lstlisting}

复制和提升的差别体现了 COW 的性能设计。若物理页仍被父子或其他映射共享，必须复制；若只剩当前映射，直接把 PTE 改回 writable 即可。无论哪种，都要处理 TLB，因为 CPU 可能还缓存着旧的只读或旧的可写翻译。

\subsection{为什么 fork 时父进程也必须写保护}

常见误解是“fork 只要把子进程页表设成只读”。这是错的。父进程如果保留 writable PTE，它可以继续写共享物理页，子进程会直接看到父进程的新内容，COW 就失效了。正确做法是父子双方都写保护，并清理父进程可能已有的 writable TLB 项。

\begin{lstlisting}
fork private pages:
  for each writable private PTE:
    increment physical page refcount
    clear parent writable bit
    mark parent PTE as COW
    install child read-only COW PTE
    invalidate parent stale writable TLB entry
\end{lstlisting}

这一步把“fork 很快”变成可理解的状态机：fork 不复制所有页，只复制页表项和引用计数；第一次写才复制物理页。代价是 fork 时要批量改 PTE 权限，并承担 TLB 失效成本。

\subsection{TLB shootdown：COW 的隐藏正确性条件}

如果父进程在 fork 前刚写过某页，CPU TLB 可能缓存了 writable 翻译。fork 清掉父 PTE writable 后，若没有失效旧 TLB，父进程下一次写可能不触发 \#PF，而是直接修改共享物理页。这不是性能问题，而是隔离正确性问题。

\begin{lstlisting}
stale TLB failure:
  parent writes page, TLB caches writable PTE
  fork clears parent PTE writable but forgets TLB invalidate
  parent writes again through stale TLB
  child observes modified shared frame
  COW handler never runs
\end{lstlisting}

多核时问题更明显：同一个地址空间的线程可能在多个 CPU 上运行，远端 CPU 也可能缓存旧翻译。内核需要本地 TLB invalidate，也可能需要 IPI 触发远端 CPU 执行 shootdown。后面看到 \code{mprotect}、\code{munmap}、COW、页迁移性能问题时，要把这笔成本算进去。

\subsection{用户指针 fault：为什么有时是 \code{-EFAULT} 而不是信号}

同一个坏地址，在用户指令里访问和在系统调用参数复制时访问，结果不同。用户程序自己执行 \code{mov rax, qword ptr [rdi]} 访问坏地址，通常应得到信号。若用户把坏指针传给 \code{read} 或 \code{write}，内核在 \code{copy\_to\_user} 或 \code{copy\_from\_user} 中 fault，则系统调用应返回 \code{-EFAULT}。

\begin{lstlisting}
; Intel syntax, user instruction fault example
mov rax, qword ptr [rdi]

; syscall usercopy fault example, conceptual
copy_from_user(kernel_buffer, user_pointer, length)
if page fault is not repairable:
  return -EFAULT
\end{lstlisting}

Linux 用异常表一类机制标记某些内核访问用户内存的指令：这里 fault 不代表内核坏了，而是用户参数不可访问，应跳到修复位置返回错误。没有这个区分，坏用户指针就会把内核打崩。

\subsection{可观察证据：不要只看 maps}

调试缺页路径时，\code{/proc/<pid>/maps} 只能告诉你 VMA，不告诉你每页是否已经有 PTE、是否驻留、是否 dirty。需要把多个证据合起来看：

\begin{longtable}{p{0.30\linewidth}p{0.30\linewidth}p{0.28\linewidth}}
\toprule
命令或文件 & 能看到什么 & 不能单独推出什么 \\
\midrule
\code{/proc/<pid>/maps} & VMA 区间、权限、文件映射 & 物理页是否已分配 \\
\code{/proc/<pid>/smaps} & RSS、dirty、shared/private 统计 & 每一次 fault 的源码分支 \\
\code{perf stat -e page-faults} & page fault 总量 & major/minor 和 fault 原因细节 \\
\code{perf stat -e minor-faults,}\\
\code{major-faults} & 缺页是否涉及 I/O & COW 与匿名零页的区别 \\
\code{strace} & 系统调用返回 \code{EFAULT} 等错误 & 用户指令自己的 page fault \\
\bottomrule
\end{longtable}

真正的分析要能把现象落回路径。例如 VMA 存在但 RSS 没涨，说明还没实际触页；minor fault 增加但 major fault 不增，通常没有存储 I/O；fork 后父子写同一页 RSS 增加，可能是 COW 分裂；系统调用返回 \code{EFAULT}，说明坏用户指针被 usercopy 路径捕获。

\subsection{配套 \cpp20 模型覆盖哪些失败}

下面的模型不是完整 Linux，也不模拟真实指令译码。它只把第一阶段最容易讲空的断言变成测试：

\begin{itemize}
\tightlist
\item 合法匿名 VMA 第一次访问会 minor fault，分配零页后重试成功。
\item 只读 VMA 的写入不能被当成 COW，必须失败且不能偷偷分配页。
\item fork 后父子共享页，第一次写触发 COW，父子内容分离。
\item 如果 fork 时忘记失效父进程 writable TLB，父进程会绕过 COW 污染子进程。
\item 文件映射第一次 fault 通过 page cache 建立映射，并记录 major fault。
\item syscall usercopy 遇到坏用户指针返回失败，而不是模拟成内核崩溃。
\item 非 canonical 地址不能被误讲成普通 VMA miss。
\end{itemize}

\begin{rubricbox}
第一阶段关于虚拟内存的最低验收：能解释一条用户态写指令为什么可能正常完成、minor fault 后完成、major fault 后完成、COW 后完成、返回 \code{-EFAULT}，或变成信号；并能指出每种结果依赖哪个 VMA/PTE/TLB 条件。
\end{rubricbox}

\subsection{完整 \cpp20 模型源码}

\lstinputlisting[language=C++,basicstyle=\ttfamily\tiny]{listings/x86_64_page_fault_cow_model.cpp}
